\section{Plan de formación y revisión}

La formación se plantea como una sesión breve seguida de práctica en el propio proyecto. El equipo puede distribuirla en tres bloques de aproximadamente 30--45 minutos, ajustando la duración a la disponibilidad y al nivel de experiencia.

\subsection{Bloque 1: interfaz accesible}
Revisar contraste, significado no dependiente del color, tipografía y uso de mayúsculas, mensajes integrados en vez de ventanas emergentes, teclado y foco. Actividad: auditar una pantalla real y proponer cambios verificables.

\subsection{Bloque 2: recorrido de backend}
Dibujar la petición, validación, permisos, almacenamiento y respuesta de una función. Actividad: analizar la creación de una tarea y enumerar casos válidos, inválidos, no autorizados y de fallo técnico.

\subsection{Bloque 3: calidad compartida}
Practicar criterios de aceptación, revisión entre compañeros, pruebas, registro de incidencias y criterio de terminado. Actividad: revisar un cambio pequeño con la lista de comprobación del manual.

\subsection{Ritmo de revisión}
Para cada funcionalidad nueva o modificada, la persona responsable de implementarla completa primero una auto-revisión. Después, otra persona del equipo revisa los criterios de aceptación y las pruebas. Julián apoya la revisión de accesibilidad y Yeray la de calidad; ambos pueden pedir cambios y registrar puntos pendientes, pero los criterios son responsabilidad de todo el equipo.

Cuando una decisión dependa del contexto del alumnado o del personal del centro, se solicitará validación al centro antes de darla por cerrada. Las pruebas automáticas o internas no sustituyen esa validación.